\chapter{Sparse Linear Systems and Krylov Subspaces}
\label{ch:krylov}

\chapterorient{Part~II handed us a linear system: discretize the weak form and
out comes $\mat{A}\vx=\vb$, with $\mat{A}$ enormous and almost entirely zeros.
Part~III handed us the machinery to run a loop without mutating anything, and the
insight that a solver is just an operator you apply. This chapter is where the two
meet. We ask why nobody solves a system this large by Gaussian elimination, why
the winning strategy instead builds the answer out of repeated \emph{applications}
of the operator---and why that is exactly the one thing the matrix-free operator of
\cref{ch:matrixfree} knows how to do. We meet the \emph{Krylov subspace}, the
nested ladder of approximations every method in Part~IV climbs; we isolate the
single \emph{step kernel} every Krylov method shares---one operator apply, a few
inner products, a few vector updates---and we read the whole solve as an
\code{iterate\_while} fold (\cref{ch:fold}) that produces a constructed operator
(\cref{ch:operators}): an approximate inverse you hold and apply. By the end you
will be able to place CG, GMRES, and their kin as specializations of one kernel,
and to say precisely which laws a solve obeys and which it deliberately does not.}

\section{The system Part~II built}
\label{sec:krylov-system}

Every analysis in this book passes through the same middle stage. Assembly
(\cref{ch:assembly,ch:galerkin}) turns the continuous weak form into a discrete
linear system
\[
  \mat{A}\,\vx \;=\; \vb,
\]
where $\vx\in\Real^N$ is the vector of unknown degree-of-freedom coefficients,
$\vb\in\Real^N$ is the discretized right-hand side, and $\mat{A}\in\Real^{N\times
N}$ is the assembled operator---the stiffness $\Kop$, the mass $\Mop$, or some
combination built from them. For the electrostatic and magnetostatic analyses
$\mat{A}$ is the symmetric positive-definite stiffness operator; for the driven
analysis it is the complex shifted system $\Kop+i\omega\Cop-\omega^2\Mop$; for a
single implicit time step of the transient analysis it is again an SPD combination.
In every case the \emph{shape} of the problem is the same: a large square system
to be solved for $\vx$.

Two features of this system govern everything that follows, and both come straight
from the finite-element construction of Part~II.

\paragraph{It is large.} $N$ is the number of degrees of freedom in the
discretization, and for a three-dimensional device meshed finely enough to
resolve the fields, $N$ runs from the hundreds of thousands into the tens of
millions. A production electromagnetic simulation routinely solves systems with
$N>10^6$. The matrix $\mat{A}$, if we wrote out all its entries, would be an
$N\times N$ array---$10^{12}$ numbers at $N=10^6$. We will see in a moment that we
never write it out.

\paragraph{It is sparse.} Here is the saving grace. The basis functions of the
finite-element method have \emph{local support}: each one is nonzero only on the
handful of mesh elements touching its node (\cref{ch:refelement}). The matrix entry
$A_{ij}=a(\phi_j,\phi_i)$ is an integral of a product of basis functions $i$ and
$j$, and that integral is zero unless the supports of $\phi_i$ and $\phi_j$
\emph{overlap}---unless nodes $i$ and $j$ share an element. Almost no two nodes do.
So almost every entry of $\mat{A}$ is exactly zero. Each row has only a small,
bounded number of nonzeros---a couple of dozen, independent of $N$---set by how
many neighbors a node has, not by how big the mesh is.

\begin{figure}[t]
\centering
\begin{tikzpicture}
\begin{axis}[
  width=0.62\linewidth, height=0.62\linewidth,
  title={\footnotesize sparsity pattern of $\mat{A}$ ($N=36$, from local support)},
  title style={yshift=-1mm},
  xmin=-0.5, xmax=35.5, ymin=-0.5, ymax=35.5,
  y dir=reverse, axis equal image,
  xtick=\empty, ytick=\empty,
  axis line style={simGray},
  xlabel={\footnotesize column $j$ (node)}, ylabel={\footnotesize row $i$ (node)},
  label style={font=\footnotesize},
]
% banded + local-support sparsity for a 6x6 grid laplacian-like stencil
\addplot[only marks, mark=square*, mark size=1.7pt, simBlue] coordinates {
  (0,0)(1,0)(6,0)
  (0,1)(1,1)(2,1)(7,1)
  (1,2)(2,2)(3,2)(8,2)
  (2,3)(3,3)(4,3)(9,3)
  (3,4)(4,4)(5,4)(10,4)
  (4,5)(5,5)(11,5)
  (0,6)(6,6)(7,6)(12,6)
  (1,7)(6,7)(7,7)(8,7)(13,7)
  (2,8)(7,8)(8,8)(9,8)(14,8)
  (3,9)(8,9)(9,9)(10,9)(15,9)
  (4,10)(9,10)(10,10)(11,10)(16,10)
  (5,11)(10,11)(11,11)(17,11)
  (6,12)(12,12)(13,12)(18,12)
  (7,13)(12,13)(13,13)(14,13)(19,13)
  (8,14)(13,14)(14,14)(15,14)(20,14)
  (9,15)(14,15)(15,15)(16,15)(21,15)
  (10,16)(15,16)(16,16)(17,16)(22,16)
  (11,17)(16,17)(17,17)(23,17)
  (12,18)(18,18)(19,18)(24,18)
  (13,19)(18,19)(19,19)(20,19)(25,19)
  (14,20)(19,20)(20,20)(21,20)(26,20)
  (15,21)(20,21)(21,21)(22,21)(27,21)
  (16,22)(21,22)(22,22)(23,22)(28,22)
  (17,23)(22,23)(23,23)(29,23)
  (18,24)(24,24)(25,24)(30,24)
  (19,25)(24,25)(25,25)(26,25)(31,25)
  (20,26)(25,26)(26,26)(27,26)(32,26)
  (21,27)(26,27)(27,27)(28,27)(33,27)
  (22,28)(27,28)(28,28)(29,28)(34,28)
  (23,29)(28,29)(29,29)(35,29)
  (24,30)(30,30)(31,30)
  (25,31)(30,31)(31,31)(32,31)
  (26,32)(31,32)(32,32)(33,32)
  (27,33)(32,33)(33,33)(34,33)
  (28,34)(33,34)(34,34)(35,34)
  (29,35)(34,35)(35,35)
};
\end{axis}
\end{tikzpicture}
\caption[Sparsity from local support]{The sparsity pattern of the assembled
operator $\mat{A}$ for a small structured mesh (each marked cell is a nonzero
entry; blanks are exact zeros). The pattern is a direct fingerprint of the mesh
connectivity: $A_{ij}\neq0$ only when basis functions $i$ and $j$ share an
element, which happens only for neighboring nodes (\cref{ch:refelement}). The
diagonal and a few near-diagonal bands carry the weight; the vast off-diagonal
plain is zero. Each row holds a bounded number of nonzeros regardless of $N$, so
the count of stored numbers grows like $O(N)$, not $O(N^2)$. This single
structural fact is what makes solving a million-unknown system possible at all.}
\label{fig:krylov-sparsity}
\end{figure}

\Cref{fig:krylov-sparsity} draws the pattern for a small structured mesh. The
nonzeros cluster on the diagonal and a few neighboring bands---the fingerprint of
the mesh's connectivity---and everything else is blank. The number of \emph{stored}
numbers is the number of nonzeros, which is $O(N)$: a constant number per row,
times $N$ rows. A million-unknown operator that would need $10^{12}$ entries dense
needs only a few times $10^7$ stored sparse. Sparsity is the only reason a problem
this large is tractable at all.

\begin{physicsbox}
The sparsity is not an accident of clever bookkeeping---it is the discrete shadow
of \emph{locality} in the physics. Maxwell's equations are local: the field at a
point is coupled directly only to the field in its immediate neighborhood, through
derivatives. The finite-element basis inherits that locality as compact support,
and the assembled operator inherits it as sparsity. A dense row in $\mat{A}$ would
mean a degree of freedom coupled directly to every other---a nonlocal physics the
weak form does not have. Sparsity is locality, written as a matrix.
\end{physicsbox}

\section{Why not just factor it?}
\label{sec:krylov-direct}

A first linear-algebra course solves $\mat{A}\vx=\vb$ by \emph{Gaussian
elimination}: reduce $\mat{A}$ to triangular form by row operations, then back-
substitute. Packaged as a reusable object this is the $LU$ factorization,
$\mat{A}=LU$ with $L$ lower- and $U$ upper-triangular; for the SPD systems of
electrostatics it specializes to the Cholesky factorization $\mat{A}=LL^{\mathsf
T}$. Once $L$ and $U$ are in hand, each solve is two cheap triangular sweeps. These
are the \emph{direct} methods: in exact arithmetic they produce the exact solution
in a finite, predetermined number of operations. Why not simply use one?

The answer is \emph{fill-in}, and it is the central fact that pushes large
finite-element solves away from direct methods. When elimination clears the entries
below a pivot, it adds multiples of the pivot row to the rows beneath. Those
additions create nonzeros where $\mat{A}$ had exact zeros: an entry $A_{ij}=0$ that
sits ``downstream'' of two nonzeros in the same elimination step becomes nonzero in
$U$. The factors $L$ and $U$ are therefore \emph{denser} than $\mat{A}$---often
dramatically so. A matrix with $O(N)$ nonzeros can have factors with
$O(N^{4/3})$, $O(N^{3/2})$, or worse nonzeros in three dimensions, and the
factorization work climbs faster still. For a 3-D mesh, sparse Cholesky typically
costs about $O(N^2)$ arithmetic and produces factors of about $O(N^{4/3})$
nonzeros---and a dense factorization, ignoring sparsity entirely, costs the
textbook $O(N^3)$.

\begin{figure}[t]
\centering
\begin{tikzpicture}
\begin{axis}[
  width=0.86\linewidth, height=62mm,
  xlabel={problem size $N$ (degrees of freedom)},
  ylabel={work (arithmetic operations)},
  xmode=log, ymode=log,
  xmin=1e3, xmax=1e7, ymin=1e3, ymax=1e15,
  xtick={1e3,1e4,1e5,1e6,1e7},
  ytick={1e3,1e6,1e9,1e12,1e15},
  legend pos=north west, legend cell align=left,
  grid=both, grid style={simGray!18},
  tick label style={font=\footnotesize}, label style={font=\small},
  legend style={font=\footnotesize},
]
% dense direct  O(N^3)
\addplot[simRust, very thick, mark=*, mark size=1.4pt, domain=1e3:1e7, samples=20]
  {x^3};
\addlegendentry{dense direct $\;O(N^3)$}
% sparse direct  ~ O(N^2)  (3-D fill-in)
\addplot[simGold, very thick, dashed, mark=triangle*, mark size=1.8pt, domain=1e3:1e7, samples=20]
  {x^2};
\addlegendentry{sparse direct $\;\sim O(N^2)$ (3-D)}
% iterative  preconditioned ~ near-linear
\addplot[simBlue, very thick, mark=square*, mark size=1.6pt, domain=1e3:1e7, samples=20]
  {60*x^1.16};
\addlegendentry{iterative (preconditioned) $\;\sim O(N^{1.16})$}
\end{axis}
\end{tikzpicture}
\caption[Direct versus iterative cost]{Work versus problem size $N$, on log--log
axes, for the three strategies. A \emph{dense} direct solve costs $O(N^3)$ and is
out of the question past a few thousand unknowns. A \emph{sparse} direct solve
exploits the sparsity of \cref{fig:krylov-sparsity} but still pays for fill-in,
costing about $O(N^2)$ in 3-D---tractable to perhaps $10^5$ unknowns, then
prohibitive. A well-\emph{preconditioned iterative} solve (Part~IV's subject)
costs close to $O(N)$ per solve: a bounded number of iterations, each doing $O(N)$
work through one sparse operator apply. The curves cross early; beyond a modest
size the iterative method is not merely faster, it is the only one that finishes.}
\label{fig:krylov-cost}
\end{figure}

\Cref{fig:krylov-cost} plots the three regimes. The dense $O(N^3)$ curve leaves
the chart almost immediately. Sparse direct, at roughly $O(N^2)$ in three
dimensions, survives to perhaps $10^5$ unknowns before fill-in makes both its
memory and its arithmetic unaffordable. The iterative curve---the subject of the
rest of Part~IV---grows close to linearly in $N$ when the iteration is well
preconditioned (\cref{ch:precond,ch:multigrid}): a bounded number of steps, each
costing one sparse operator apply at $O(N)$. The crossover comes early, and on the
far side of it the iterative method is not just the faster choice; for a
ten-million-unknown 3-D solve it is the \emph{only} choice that fits in memory and
finishes in time.

\begin{notationbox}
Direct methods are not banished from the simulator. They reappear, deliberately,
as \emph{components inside} iterative solves: a sparse direct factorization makes an
excellent preconditioner on a coarse problem (\cref{ch:precond}), and the coarsest
level of a multigrid hierarchy is often solved directly (\cref{ch:multigrid}). The
claim of this section is narrower and sharper: do not solve the \emph{full}
fine-scale system directly. Use a direct method only where the problem has been
made small enough that fill-in no longer bites.
\end{notationbox}

\subsection{What an iterative method asks of the operator}
\label{sec:krylov-action}

There is a second reason the iterative route wins, and it is the one that ties this
chapter back to \cref{ch:matrixfree}. A direct method must \emph{read the entries}
of $\mat{A}$: elimination operates on $A_{ij}$ one at a time. An iterative method
does not. Look ahead to any solver in Part~IV---conjugate gradients
(\cref{ch:cg}), GMRES (\cref{ch:gmres})---and watch what it does with $\mat{A}$. It
forms inner products of vectors, it takes linear combinations of vectors, it
checks norms. The \emph{only} way it ever touches $\mat{A}$ is through one
operation: given a vector $\vv$, produce the vector $\mat{A}\vv$. It never asks for
an entry $A_{ij}$, never asks for a row or a column, never inspects the sparsity
pattern. It needs the operator's \emph{action}, and nothing else:
\[
  \vv \;\longmapsto\; \mat{A}\vv .
\]

This is precisely the interface \cref{ch:matrixfree} built. The matrix-free
operator never stores $\mat{A}$; it computes $\mat{A}\vv$ on demand as a short
chain of element-local tensor contractions---gather, evaluate, weight, integrate,
scatter. An iterative solver is the perfect consumer of that interface: it asks for
$\mat{A}\vv$, the matrix-free operator computes it, and neither ever needs the
assembled array to exist. The two ideas were built for each other. A direct method
\emph{cannot} use a matrix-free operator---it needs the entries elimination
operates on---which is one more reason direct methods do not scale to the
matrix-free, high-order, accelerator-bound regime this book targets.

\begin{keyidea}
An iterative solver needs only the operator's \emph{action}, $\vv\mapsto\mat{A}\vv$
(and, with preconditioning, $\vr\mapsto\mat{M}^{-1}\vr$)---never the entries of
$\mat{A}$. This is the licence for the whole of Part~IV. It means the solver
composes directly with the matrix-free operator of \cref{ch:matrixfree}, which
supplies exactly that action and never forms the matrix; and it is why the cost of
a solve is counted in \emph{operator applies}, the one expensive thing each step
does. \emph{One apply per step} is the cost invariant the rest of the chapter makes
precise.
\end{keyidea}

\section{The Krylov subspace}
\label{sec:krylov-subspace}

If all we may do is apply $\mat{A}$ to vectors, what can we build? Start from the
right-hand side $\vb$ (or, more generally, from the initial residual
$\vr_0=\vb-\mat{A}\vx_0$). Apply $\mat{A}$ and you get $\mat{A}\vb$. Apply again:
$\mat{A}^2\vb$. Each application is one matrix-free action; after $m-1$ of them you
hold the family
\[
  \vb,\quad \mat{A}\vb,\quad \mat{A}^2\vb,\quad \dots,\quad \mat{A}^{m-1}\vb .
\]
The space these vectors span is the \emph{Krylov subspace} of dimension $m$.

\begin{definition}[Krylov subspace]
\label{def:krylov}
For a square operator $\mat{A}\in\Real^{N\times N}$ and a nonzero vector
$\vb\in\Real^N$, the $m$-th \emph{Krylov subspace} is
\[
  \mathcal{K}_m(\mat{A},\vb)
   \;=\; \operatorname{span}\bigl\{\,\vb,\ \mat{A}\vb,\ \mat{A}^2\vb,\ \dots,\
   \mat{A}^{m-1}\vb\,\bigr\}
   \;\subseteq\; \Real^N .
\]
The spanning vectors $\vb,\mat{A}\vb,\dots$ are the \emph{Krylov sequence}; each
new one costs exactly one application of the operator's action.
\end{definition}

A Krylov method builds its approximate solution \emph{out of} this subspace. At
step $m$ it picks the vector $\vx_m\in\vx_0+\mathcal{K}_m(\mat{A},\vr_0)$ that is
``best'' by some criterion the particular method chooses---minimizing the residual
norm (GMRES, \cref{ch:gmres}), minimizing the $\mat{A}$-norm of the error (CG,
\cref{ch:cg}). The subspaces nest,
\[
  \mathcal{K}_1 \subseteq \mathcal{K}_2 \subseteq \mathcal{K}_3 \subseteq \cdots,
\]
because each is the previous one plus one more power of $\mat{A}$. So the
approximation can only improve as $m$ grows: a best-in-$\mathcal{K}_{m+1}$ vector
is at least as good as the best-in-$\mathcal{K}_m$ vector, since the smaller
subspace sits inside the larger. The method climbs this ladder of nested
subspaces, one operator apply per rung, until the residual is small enough to stop.

\begin{figure}[t]
\centering
\begin{tikzpicture}[scale=1.0]
  % nested subspaces as nested ellipses, with the true solution as a star
  \draw[draw=simBlue, fill=simBgBlue, thick] (0,0) ellipse (4.1 and 2.5);
  \draw[draw=simTeal, fill=simBgGreen, thick] (-0.5,-0.1) ellipse (2.9 and 1.75);
  \draw[draw=simGold, fill=simBgGold, thick] (-1.0,-0.15) ellipse (1.7 and 1.05);
  % labels for the subspaces
  \node[font=\scriptsize\bfseries, simGold!70!black] at (-1.0,-0.15) {$\mathcal{K}_1$};
  \node[font=\scriptsize\bfseries, simTeal] at (1.35,0.95) {$\mathcal{K}_2$};
  \node[font=\scriptsize\bfseries, simBlue] at (3.0,1.85) {$\mathcal{K}_3$};
  \node[font=\scriptsize, simGray] at (4.55,2.3) {$\Real^N$};
  % the true solution
  \node[star, star points=5, star point ratio=2.3, draw=simRust, fill=simRust,
        inner sep=1.3pt] (sol) at (2.0,-1.4) {};
  \node[font=\scriptsize\bfseries, simRust, right=1mm of sol] {$\vx^\star$ (true solution)};
  % the iterates marching toward it
  \node[circle, fill=simInk, inner sep=1.2pt] (x1) at (-1.0,-0.15) {};
  \node[circle, fill=simInk, inner sep=1.2pt] (x2) at (0.5,-0.55) {};
  \node[circle, fill=simInk, inner sep=1.2pt] (x3) at (1.5,-1.05) {};
  \node[font=\scriptsize, simInk, above left=0mm and 0mm of x1] {$\vx_1$};
  \node[font=\scriptsize, simInk, above=0mm of x2] {$\vx_2$};
  \node[font=\scriptsize, simInk, above right=0mm and 0mm of x3] {$\vx_3$};
  \draw[flow, shorten >=2pt, shorten <=2pt] (x1)--(x2);
  \draw[flow, shorten >=2pt, shorten <=2pt] (x2)--(x3);
  \draw[flow, densely dashed, simRust, shorten >=2pt, shorten <=2pt] (x3)--(sol);
\end{tikzpicture}
\caption[The nested Krylov subspaces]{The Krylov subspaces nest,
$\mathcal{K}_1\subseteq\mathcal{K}_2\subseteq\mathcal{K}_3\subseteq\cdots$, each
adding one more power of $\mat{A}$ to the span and so growing by (at most) one
dimension per step. Each costs one operator apply to extend. The method's iterate
$\vx_m$ is the best approximation to the true solution $\vx^\star$ available within
$\vx_0+\mathcal{K}_m$; because the subspaces grow, the iterates can only improve,
marching toward $\vx^\star$ as the ladder climbs. In exact arithmetic the ladder
reaches $\vx^\star$ exactly once $\mathcal{K}_m$ is large enough to contain
it---at the latest when $m$ equals the degree of $\mat{A}$'s minimal polynomial
(\cref{sec:krylov-poly}).}
\label{fig:krylov-nested}
\end{figure}

\Cref{fig:krylov-nested} draws the picture: nested subspaces, and inside them a
sequence of iterates marching toward the true solution $\vx^\star$. The whole of
Part~IV is variations on how to choose the best iterate in each subspace and how to
do it cheaply; this section is about why the subspace contains a good answer at all.

\subsection{Why powers of \texorpdfstring{$\mat{A}$}{A} capture the solution}
\label{sec:krylov-poly}

It is not obvious that repeatedly applying $\mat{A}$ should march toward
$\mat{A}^{-1}\vb$---applying $\mat{A}$ looks like the opposite of inverting it. The
reason it works is a fact about polynomials in $\mat{A}$, and it is worth seeing
because it explains both why Krylov methods converge \emph{and} how fast.

Notice first that any vector in $\mathcal{K}_m(\mat{A},\vb)$ is, by definition, a
linear combination of $\vb,\mat{A}\vb,\dots,\mat{A}^{m-1}\vb$---that is, it is
$q(\mat{A})\,\vb$ for some polynomial $q$ of degree at most $m-1$:
\[
  \vx_m \in \mathcal{K}_m
  \quad\Longleftrightarrow\quad
  \vx_m = q(\mat{A})\,\vb
  \ \text{ for some } q \text{ with } \deg q \le m-1 .
\]
So ``find the best iterate in $\mathcal{K}_m$'' is the same as ``find the best
degree-$(m-1)$ polynomial $q$ such that $q(\mat{A})\vb\approx\mat{A}^{-1}\vb$.''
Krylov methods are, under the hood, polynomial approximations of the inverse.

Now the key fact. Every square matrix satisfies a \emph{minimal polynomial}: there
is a polynomial $p$ of some degree $d\le N$ with $p(\mat{A})=0$. Write
$p(t)=c_0+c_1t+\dots+c_d t^d$ with $c_0\neq0$ (true when $\mat{A}$ is invertible).
From $p(\mat{A})=0$ we can solve for the identity:
\[
  c_0 I = -\bigl(c_1\mat{A}+\dots+c_d\mat{A}^d\bigr)
  \;\Longrightarrow\;
  I = -\tfrac{1}{c_0}\bigl(c_1\mat{A}+\dots+c_d\mat{A}^d\bigr),
\]
and multiplying by $\mat{A}^{-1}$ gives the inverse as a \emph{polynomial in
$\mat{A}$ of degree $d-1$}:
\[
  \mat{A}^{-1} \;=\; -\tfrac{1}{c_0}\bigl(c_1 I + c_2\mat{A} + \dots + c_d\mat{A}^{d-1}\bigr).
\]
Therefore $\mat{A}^{-1}\vb=-\tfrac{1}{c_0}(c_1\vb+c_2\mat{A}\vb+\dots+c_d\mat{A}^{d-1}\vb)$
is exactly a vector in $\mathcal{K}_d(\mat{A},\vb)$. The true solution \emph{lives
inside} a Krylov subspace of dimension $d\le N$. A Krylov method that picks the
best iterate in each $\mathcal{K}_m$ therefore reaches the exact solution by step
$d$ at the latest---in exact arithmetic, a finite procedure, no worse than $N$
steps.

\begin{keyidea}
A vector in $\mathcal{K}_m(\mat{A},\vb)$ is exactly $q(\mat{A})\vb$ for a polynomial
$q$ of degree $<m$, so a Krylov method is a polynomial approximation of
$\mat{A}^{-1}$. The minimal polynomial guarantees $\mat{A}^{-1}\vb$ itself lies in
$\mathcal{K}_d$ for some $d\le N$: \emph{in exact arithmetic a Krylov method
terminates with the exact solution in at most $N$ steps}. But $N$ steps would be no
bargain---the method is useful only because it gets \emph{close} far sooner, and how
soon is governed by the spectrum of $\mat{A}$ (\cref{sec:krylov-convergence}).
\end{keyidea}

\subsection{Conditioning and the speed of convergence}
\label{sec:krylov-convergence}

The polynomial picture also tells us how \emph{fast} a Krylov method converges,
and the answer is the bridge to preconditioning. The error after $m$ steps is
controlled by how small a degree-$m$ polynomial $q$ (with $q(0)=1$, the
normalization the residual forces) can be made on the \emph{eigenvalues} of
$\mat{A}$. Roughly: the residual at step $m$ is bounded by
\[
  \frac{\norm{\vr_m}}{\norm{\vr_0}}
   \;\le\; \min_{\substack{q,\ \deg q\le m \\ q(0)=1}}\ \max_{\lambda\in\sigma(\mat{A})}\ \abs{q(\lambda)},
\]
where $\sigma(\mat{A})$ is the set of eigenvalues. Convergence is fast exactly when
a low-degree polynomial, pinned to $1$ at the origin, can be made small on the
whole spectrum.

The consequence is the single most important rule of thumb in iterative solving:
\emph{a clustered spectrum converges fast; a spread-out spectrum converges slowly.}
If all the eigenvalues huddle in a tight cluster away from zero, a low-degree
polynomial can dip small across the whole cluster, and a few steps suffice. If the
eigenvalues are spread over a wide range---if the \emph{condition number}
$\kappa=\lambda_{\max}/\lambda_{\min}$ is large---then no low-degree polynomial can
be small everywhere at once, and convergence crawls. For SPD systems the classic
bound makes this quantitative: CG reduces the error by a factor that improves like
$1-2/\sqrt{\kappa}$ per step, so the iteration count scales like $\sqrt{\kappa}$.

\begin{figure}[t]
\centering
\begin{tikzpicture}
\begin{axis}[
  width=0.86\linewidth, height=60mm,
  xlabel={iteration $m$},
  ylabel={relative residual $\norm{\vr_m}/\norm{\vr_0}$},
  ymode=log,
  xmin=0, xmax=60, ymin=1e-10, ymax=2,
  xtick={0,10,20,30,40,50,60},
  ytick={1e0,1e-2,1e-4,1e-6,1e-8,1e-10},
  legend pos=north east, legend cell align=left,
  grid=both, grid style={simGray!18},
  tick label style={font=\footnotesize}, label style={font=\small},
  legend style={font=\footnotesize},
]
% clustered spectrum: fast linear-on-log decay
\addplot[simBlue, very thick, mark=square*, mark size=1.3pt, domain=0:18, samples=19]
  {exp(-1.15*x)};
\addlegendentry{clustered spectrum (small $\kappa$)}
% spread spectrum: slow decay
\addplot[simRust, very thick, mark=*, mark size=1.3pt, domain=0:60, samples=31]
  {exp(-0.35*x)};
\addlegendentry{spread spectrum (large $\kappa$)}
% tolerance line
\addplot[simGray, densely dashed, domain=0:60, samples=2] {1e-8};
\addlegendentry{tolerance}
\end{axis}
\end{tikzpicture}
\caption[Convergence depends on the spectrum]{The relative residual versus
iteration count for two systems of the same size, on a log scale (a straight line
is geometric decay). When the eigenvalues of $\mat{A}$ are \emph{clustered}
(condition number small), a low-degree polynomial is small across the whole
spectrum and the residual plunges---tolerance reached in a handful of steps.
When the spectrum is \emph{spread} (condition number large), no low-degree
polynomial can be small everywhere at once and the residual crawls down, taking
many times as many steps to cross the same tolerance. Preconditioning
(\cref{ch:precond}) is precisely the art of replacing $\mat{A}$ with an
equivalent system whose spectrum is clustered---of moving a problem from the slow
curve onto the fast one.}
\label{fig:krylov-convergence}
\end{figure}

\Cref{fig:krylov-convergence} shows the two regimes: the same tolerance, reached in
a handful of steps on a clustered spectrum and only after a long crawl on a spread
one. This is why the bare Krylov method of this chapter is never the whole story.
The method itself is fixed; what we can change is the \emph{spectrum} of the system
it sees. \emph{Preconditioning} (\cref{ch:precond,ch:multigrid}) replaces
$\mat{A}\vx=\vb$ with an equivalent system $\mat{M}^{-1}\mat{A}\vx=\mat{M}^{-1}\vb$
whose operator $\mat{M}^{-1}\mat{A}$ has a clustered spectrum---moving the problem
from the slow curve onto the fast one. The Krylov method supplies the iteration;
the preconditioner supplies the conditioning; and together they reach close to
$O(N)$ work. That is the whole arc of Part~IV, and this chapter is its foundation.

\section{The shared step kernel}
\label{sec:krylov-kernel}

We turn now from \emph{why} Krylov methods work to \emph{what} they all do, step by
step. Here is the observation that organizes the rest of Part~IV: every Krylov
method---CG, GMRES, FGMRES, MINRES, BiCGStab---has an inner step of the \emph{same
shape}. Strip away the method-specific algebra and the per-step kernel is always
the same five-part sequence:
\begin{enumerate}[nosep]
  \item \textbf{one operator apply}: form $\mat{A}\vp$ (the matrix-free action,
  \cref{ch:matrixfree})---and, if preconditioned, one preconditioner apply
  $\mat{M}^{-1}\vr$;
  \item \textbf{a few inner products and norms}: $\inner{\cdot}{\cdot}$ and
  $\norm{\cdot}$ to form the step's scalar coefficients and measure the residual;
  \item \textbf{a few vector updates} (\emph{axpy}): $\vy\leftarrow\alpha\vx+\vy$,
  combining vectors with the freshly computed scalars;
  \item \textbf{an orthogonalization} (for the basis-building methods): make the new
  direction orthogonal to the previous ones (\cref{ch:orthog});
  \item \textbf{a state update}: fold the new iterate, residual, and scalars into
  the carry, and record the residual norm for the convergence test.
\end{enumerate}
We call this kernel \code{krylov\_step}. A particular method---CG, GMRES---is a
particular \emph{specialization} of it: a choice of which scalars, which updates,
whether the orthogonalization stage runs. But the skeleton is fixed, and one fact
about it is load-bearing.

\begin{keyidea}
Every Krylov method shares one step kernel, \code{krylov\_step}: one operator apply,
a few inner products and norms, a few \emph{axpy} updates, an optional
orthogonalization, and a state update. A specific method is a specialization---it
binds the scalars and the updates---but the shape is invariant. The \emph{cost
invariant} is the heart of it: \emph{exactly one operator apply per step}. Since the
apply is the one expensive operation (it touches all $N$ unknowns through the sparse
operator; everything else is cheap vector arithmetic), counting applies counts the
cost of the solve, and ``one apply per step'' is the budget every Krylov method
spends.
\end{keyidea}

\begin{figure}[t]
\centering
\begin{tikzpicture}[node distance=7mm and 10mm]
  % input carry
  \node[data, text width=20mm] (in) {carry $K$\\\footnotesize $\vx,\vr,\vp,\dots$};
  % stage 1: operator apply (the one expensive op)
  \node[stage, right=of in, text width=20mm, fill=simBgRust, draw=simRust] (apply)
    {\footnotesize \textbf{apply}\\$\vw=\mat{A}\vp$\\\scriptsize the one $O(N)$ op};
  % stage 2: inner products
  \node[opbox, right=of apply, text width=20mm] (dots)
    {\footnotesize \textbf{dots/norms}\\$\inner{\vp}{\vw},\ \norm{\vr}$};
  % stage 3: orthogonalization (optional)
  \node[extern, right=of dots, text width=20mm] (orth)
    {\footnotesize \textbf{orthog.}\\\scriptsize (GMRES only)};
  % stage 4: axpy updates
  \node[opbox, right=of orth, text width=20mm] (axpy)
    {\footnotesize \textbf{updates}\\\scriptsize \emph{axpy}: $\vx,\vr,\vp$};
  % output carry
  \node[data, right=of axpy, text width=20mm] (out) {carry $K'$\\\footnotesize $+$ residual};
  \draw[flow] (in) -- (apply);
  \draw[flow] (apply) -- node[above,font=\scriptsize]{$\vw$} (dots);
  \draw[flow] (dots) -- node[above,font=\scriptsize]{scalars} (orth);
  \draw[flow] (orth) -- (axpy);
  \draw[flow] (axpy) -- (out);
  % residual peels off downward for the convergence test
  \node[product, below=10mm of axpy, text width=24mm] (res)
    {\footnotesize residual norm\\\scriptsize $\to$ convergence test};
  \draw[refedge] (axpy.south) -- (res.north);
  % the cost-invariant brace under the apply
  \draw[decorate, decoration={brace, amplitude=4pt, mirror}, simRust]
    ([yshift=-2mm]apply.south west) -- ([yshift=-2mm]apply.south east)
    node[midway, below=4pt, font=\scriptsize, simRust, align=center]
    {one apply\\per step};
\end{tikzpicture}
\caption[The shared Krylov step kernel]{The kernel \code{krylov\_step} every Krylov
method shares, as a dataflow. A carry $K$ (the working vectors $\vx,\vr,\vp$ and
scalars) enters; the \textbf{one operator apply} $\vw=\mat{A}\vp$ (rust, the single
$O(N)$ operation) runs once; a few \textbf{inner products and norms} form the step's
scalars; the basis-building methods run an \textbf{orthogonalization}
(\cref{ch:orthog}; absent in CG, dashed); a few \emph{axpy} \textbf{updates} combine
vectors with the new scalars into the next carry $K'$; and the residual norm peels
off for the convergence test. The method (CG, GMRES, \dots) chooses the scalars and
which updates run; the \emph{shape}---and the one-apply-per-step cost---is fixed.}
\label{fig:krylov-kernel}
\end{figure}

\Cref{fig:krylov-kernel} draws the kernel as a dataflow. Note what is and is not
expensive. The operator apply touches all $N$ unknowns through the sparse operator;
it is the $O(N)$ operation, and it is the only one. The inner products and norms are
$O(N)$ reductions but with tiny constants; the \emph{axpy} updates are $O(N)$
elementwise; the orthogonalization, when present, costs a handful of inner products
against the existing basis. Everything but the apply is cheap vector arithmetic. So
the apply is the currency in which we price a solve, and the kernel spends exactly
one unit of it per step. When the apply is matrix-free (\cref{ch:matrixfree}), that
one apply is itself a chain of element-local contractions---but it is still one
apply, and the cost invariant holds.

\section{The solver as a fold, and as an operator}
\label{sec:krylov-fold}

We have the kernel; now we wrap a loop around it. The outer structure of a Krylov
solve is exactly the \code{iterate\_while} fold of \cref{ch:fold}: carry the
solver's working state forward, take one \code{krylov\_step} per pass, test a
convergence predicate, stop when the residual is small enough. There is nothing new
to invent---the solver is a textbook instance of the combinator the framework
already built.

Recall the kernel signature from the framework: the value-threaded Krylov step
maps a carry to the next carry plus a demand-prunable readout,
\begin{lstlisting}[style=l4style]
-- one step: thread the working bundle K, the persistent state s; report extras
krylov_step :: (op, K, s) -> (K', s', outputs)
\end{lstlisting}
where \code{op} is the construction-bound operator data (the system operator and
preconditioner, frozen in; \cref{ch:operators}), \code{K} is the iterate bundle
(the working vectors and scalars), \code{s} is the persistent solver state (the
iteration counter, the converged flag, the residual proxy), and \code{outputs} is
the per-step readout (typically the residual norm) that the trajectory collects.
Folding this kernel with \code{iterate\_while} is the entire outer loop:
\begin{lstlisting}[style=l4style]
-- the Krylov solve IS a fold of the step kernel (schematically)
solve op K0 s0 =
  iterate_while
    { K: K0, s: s0 }                              -- initial carry
    (\c -> not c.s.converged && c.s.it < max_it)  -- predicate: pure on the carry
    (\c -> let (K', s', outs) = krylov_step op c.K c.s
           in { state: { K: K', s: s' }, residual_norm: outs.residual_norm })
\end{lstlisting}
The predicate reads only the carry---the converged flag and iteration count it
carries---exactly as the predicate rule of \cref{sec:predicaterule} demands; the
residual the step computes is folded into the carry's \code{s} so the \emph{next}
pass's predicate can read it. The per-step residual norm rides out in the extras and
into the trajectory, where a convergence monitor reads it when asked---and, by
demand-driven pruning (\cref{sec:pruning}), is not computed at all when nothing
reads it. The solver is a fold, and it inherits every property the fold has.

\begin{figure}[t]
\centering
\begin{tikzpicture}[font=\footnotesize, node distance=6mm]
  % carry spine threading SimState + Krylov bundle
  \node[data, minimum width=16mm, align=center] (c0) at (0,0) {$K_0,s_0$};
  \node[diamond, draw=simBlue, fill=simBgBlue, aspect=1.7, inner sep=1pt,
        font=\scriptsize] (p0) at (2.1,0) {conv.?};
  \node[opbox, minimum width=18mm] (st0) at (4.3,0) {\code{krylov\_step}};
  \node[data, minimum width=16mm, align=center] (c1) at (6.6,0) {$K_1,s_1$};
  \node[diamond, draw=simBlue, fill=simBgBlue, aspect=1.7, inner sep=1pt,
        font=\scriptsize] (p1) at (8.7,0) {conv.?};
  \node[opbox, minimum width=18mm] (st1) at (10.9,0) {\code{krylov\_step}};
  \node[font=\scriptsize, simGray] (more) at (12.6,0) {$\cdots$};
  \draw[flow] (c0)--(p0);
  \draw[flow] (p0)-- node[above,font=\scriptsize]{no} (st0);
  \draw[flow] (st0)-- node[above,font=\scriptsize]{\code{.state}} (c1);
  \draw[flow] (c1)--(p1);
  \draw[flow] (p1)-- node[above,font=\scriptsize]{no} (st1);
  \draw[flow] (st1)-- (more);
  % converged exits down to the result
  \node[product, minimum width=22mm] (fin) at (8.7,-2.2)
    {\code{SolveResult}: $\vx$, converged};
  \draw[backflow] (p0.south) |- (fin.west);
  \draw[backflow] (p1.south) -- (fin.north);
  % residual trajectory peels downward
  \node[data, fill=simBgGold, draw=simGold, minimum width=40mm] (traj) at (7.7,-1.2)
    {residual history $[\,r_0,\,r_1,\,\dots\,]$ (pruned if unread)};
  \draw[refedge] (st0.south) -- (st0.south|-traj.north);
  \draw[refedge] (st1.south) -- (st1.south|-traj.north);
\end{tikzpicture}
\caption[A Krylov solve as an \texttt{iterate\_while} fold]{The outer loop of a
Krylov solve is the \code{iterate\_while} fold of \cref{ch:fold}. The carry
threads the solver's state---the iterate bundle $K$ (working vectors and scalars)
and the persistent state $s$ (iteration counter, converged flag, residual
proxy)---forward along the top. Before each pass the predicate tests convergence
on the \emph{current carry} (the \texttt{while} convention, \cref{sec:predicaterule});
on ``not converged'' a single \code{krylov\_step} runs, emitting the next carry. The
per-step residual peels downward into the trajectory, where a monitor reads it---or,
unread, it is pruned (\cref{sec:pruning}). The first ``converged'' exits the loop,
surfacing the iterate $\vx$ as the \code{SolveResult}. Compare \cref{fig:threaded}:
this is that figure with ``step'' the Krylov kernel and the carry the solver state.}
\label{fig:krylov-iterwhile}
\end{figure}

\Cref{fig:krylov-iterwhile} draws the fold for the Krylov solve specifically---the
threaded-carry picture of \cref{ch:fold}, with \code{krylov\_step} in the step box
and the residual history peeling off as the trajectory. The state-threaded picture
also explains why the carry splits in two (\cref{ch:strata}): the persistent
state $s$ (counter, flags, the externally visible iterate) is one stratum, and the
ephemeral iterate bundle $K$ (the workspace vectors, reborn at each restart of a
restarted method) is another. The fold threads both; their different lifetimes are
why GMRES's restart logic wraps an \emph{outer} loop around the inner fold rather
than flattening the two (a fold-merge the non-laws of \cref{sec:laws} forbid, and
which \cref{ch:gmres} respects).

\subsection{The whole solve is a constructed operator}
\label{sec:krylov-operator}

Step back from the loop and look at the solve as a single object. It takes a
right-hand side $\vb$ and returns an approximate solution $\vx\approx\mat{A}^{-1}\vb$.
That is the signature of an \emph{operator}: vector in, vector out. The entire
Krylov solve is a \emph{constructed operator} in the exact sense of
\cref{ch:operators}---built once with its system operator, preconditioner,
tolerances, and iteration budget baked into immutable internal state, then applied to
a right-hand side as a pure function. We met this as \code{ksp\_solve}; here is its
signature.

\begin{lstlisting}[style=l4style]
-- the solver is a constructed operator: an approximate inverse you apply
ksp_solve :: (K: Solver[A], b: Tensor[N]) -> SolveResult[N]

SolveResult[N] = {
  x          : Tensor[N],   -- approximate solution to A . x = b
  converged  : Bool,        -- did the convergence test pass?
  iterations : Int,         -- inner Krylov iterations consumed
  initial_res: Real,        -- initial residual norm
  final_res  : Real         -- final residual norm
}
\end{lstlisting}

The construction-bound value \code{K} carries the system operator \code{A}, an
optional preconditioner, and the convergence-control parameters; \emph{all} the
per-method choices---CG versus GMRES versus FGMRES, preconditioner side,
orthogonalization method---are absorbed into \code{K} at construction
(\cref{ch:operators}), so the per-call surface $(K,\vb)$ is variant-free. The solve
is the constructed-operator pattern of \cref{ch:operators} in its most consequential
instance: $\code{ksp\_solve}(K,\cdot)$ is an \emph{approximate inverse} you hold and
apply, indistinguishable in form from applying $\mat{A}$ itself. This is exactly the
\emph{solver-as-operator} move that lets one GMRES drive any preconditioner
(\cref{ch:operators}): the preconditioner is just another \code{apply}-shaped value
the kernel calls.

\begin{figure}[t]
\centering
\begin{tikzpicture}[node distance=10mm]
  % construction inputs (build-time)
  \node[data, align=left, text width=26mm] (cfg)
    {operator $\mat{A}$\\preconditioner\\tolerances,\\iteration cap};
  \node[stage, minimum width=24mm, minimum height=14mm, right=12mm of cfg] (con)
    {\textbf{construct}\\[1pt]\footnotesize freeze config,\\absorb method};
  \draw[flow] (cfg) -- node[above,font=\scriptsize]{once} (con);
  % the opaque solver-operator value
  \node[draw=simBlue, fill=simBgBlue, very thick, rounded corners=3pt,
        minimum width=26mm, minimum height=20mm, right=12mm of con] (op) {};
  \node[font=\ttfamily\scriptsize, simInk] at (op.north) [yshift=-3.4mm]
    {ksp\_solve};
  \node[font=\scriptsize\itshape, simGray, align=center] at (op.center)
    {$\approx\mat{A}^{-1}$\\(frozen:\\$\mat{A}$, pc,\\tols)};
  \draw[flow] (con) -- (op);
  % many applies out
  \foreach \y/\lbl in {1.0/a, 0.0/b, -1.0/c}{
    \node[draw=simTeal, fill=simBgGreen, semithick, rounded corners=2pt,
          minimum width=20mm, minimum height=6mm] (ap\lbl) at ($(op.east)+(2.9,\y)$)
          {\footnotesize\ttfamily ksp\_solve(K, b)};
    \draw[flow] (op.east) -- (ap\lbl.west);
  }
  \node[font=\scriptsize\itshape, simGray] at ($(op.east)+(2.9,-1.7)$)
    {many RHS, each a fold inside};
  % band labels
  \node[font=\scriptsize\bfseries, simRust] at ($(con)+(0,-1.5)$) {build-time};
  \node[font=\scriptsize\bfseries, simGreen] at ($(op.east)+(2.9,1.7)$) {run-time};
\end{tikzpicture}
\caption[The solve as a constructed operator]{The whole Krylov solve is a
constructed operator (\cref{ch:operators}). \emph{Construction} (left, build-time,
once) freezes the system operator $\mat{A}$, the preconditioner, and the
convergence-control parameters into an opaque value \code{K}, absorbing the choice
of method along the way. \emph{Application} (right, run-time, many times) feeds a
right-hand side $\vb$ and runs the \code{iterate\_while} fold of
\cref{fig:krylov-iterwhile} inside, returning $\vx\approx\mat{A}^{-1}\vb$. The
solve wears the same \code{apply} interface as the operator it inverts: it is an
\emph{approximate inverse you apply}, and the fixed-operator map of
\cref{ch:situating} (electrostatics solving against many right-hand sides) is
exactly this operator applied repeatedly with $\mat{A}$ hoisted out of the loop.}
\label{fig:krylov-ksp}
\end{figure}

\Cref{fig:krylov-ksp} draws it: build once, apply many. This is also the resolution
of the fixed-operator map from \cref{ch:situating}. Electrostatics solves
$\Kop\vx_i=\vb_i$ for many right-hand sides against \emph{one} operator; in this
language, it constructs \code{ksp\_solve} once with $\Kop$ frozen in and applies it to
each $\vb_i$, hoisting the expensive construction out of the loop. The constructed
operator \emph{is} the hoist.

\subsection{What the solve guarantees---and what it deliberately does not}
\label{sec:krylov-laws}

Because the solve is an approximate inverse, it obeys the laws an inverse obeys---in
the limit of a tight tolerance---and, just as importantly, it pointedly fails to
obey several laws an \emph{exact} inverse would. Knowing both is what keeps a caller
from trusting a guarantee the solve does not make. The laws below hold modulo the
convergence tolerance (equalities in the limit $\text{rel\_tol}\to0$,
$\text{max\_it}\to\infty$, treating $K$ as the formal map $\vb\mapsto\mat{A}^{-1}\vb$).

\begin{description}[style=nextline, leftmargin=1.4em]
  \item[Linearity in the right-hand side.] $\code{ksp\_solve}(K,\alpha\vb_1+\beta\vb_2).\vx
  \approx \alpha\,\code{ksp\_solve}(K,\vb_1).\vx + \beta\,\code{ksp\_solve}(K,\vb_2).\vx$.
  This is the linearity of $\mat{A}^{-1}$, exact in the formal limit and approximate
  at finite tolerance (the gap bounded by the condition number). The
  \emph{statistics} fields are not linear: different right-hand sides take different
  iteration counts and produce different residual histories.
  \item[Near-inverse on the right-hand side.] $\code{ksp\_solve}(K,\mat{A}\vx).\vx
  \approx \vx$: applying the solve to $\mat{A}\vx$ recovers $\vx$ to within the
  tolerance. This is the defining property---the solve is an approximate
  $\mat{A}^{-1}$.
  \item[Zero right-hand side gives zero.] $\code{ksp\_solve}(K,\mathbf{0}).\vx=\mathbf{0}$,
  exactly, not merely modulo tolerance (the iteration short-circuits on a zero
  initial residual).
\end{description}

Now the non-laws---the guarantees the solve refuses to make, every one of them
load-bearing in the sense of \cref{ch:bigidea}.

\begin{pitfall}
A Krylov solve is \emph{not} a bit-exact, deterministic function of $(K,\vb)$. Three
distinct sources of non-determinism are deliberate and must not be ``optimized
away.''
\begin{itemize}[nosep]
  \item \textbf{Reduction-order non-determinism.} The inner products and norms of the
  kernel are floating-point reductions, and floating-point addition is not
  associative. Reordering a reduction tree (as a parallel backend will) changes the
  bit-level scalars, hence the bit-level iterate count and the bit-level $\vx$. The
  mathematical solution is unchanged; its floating-point realization is not.
  \item \textbf{Orthogonalization-variant non-determinism.} For the basis-building
  methods, the orthogonalization choice (modified versus classical Gram--Schmidt,
  \cref{ch:orthog}) changes the per-step arithmetic and so the bit-level trajectory.
  Same answer, different bits.
  \item \textbf{Initial-guess non-determinism.} A warm start and a cold start take
  different paths through the Krylov subspaces and converge to bit-different $\vx$
  in different iteration counts. The mathematical solution is the same.
\end{itemize}
And the solve does \emph{not} compose exactly with the operator: $\mat{A}\,
\code{ksp\_solve}(K,\vb).\vx=\vb$ holds only \emph{approximately}, within the
tolerance, never exactly at finite precision. An algorithm that assumes the residual
is exactly zero after a solve---some iterative-refinement schemes do---must guard.
\emph{Convergence depends on conditioning, not on exactness}; the solve is a
controlled approximation, and its non-determinism is the honest price of running it
on a parallel, finite-precision machine.
\end{pitfall}

These non-laws are not defects to be apologized for; they are the truthful contract.
A Krylov solve is a controlled approximation whose accuracy you dial with a
tolerance and whose speed you buy with conditioning. Pretending it is an exact,
deterministic inverse would be convenient and false---and would silently corrupt any
caller that leaned on the pretense.

\section{Worked examples}
\label{sec:krylov-wex}

\begin{workedexample}{Building $\mathcal{K}_3(\mat{A},\vb)$ for a $3\times3$ SPD
system}{krylovsub}
We build a Krylov subspace explicitly and watch the solution appear inside it. Take
the small SPD matrix and right-hand side
\[
  \mat{A}=\begin{pmatrix}2&1&0\\1&2&1\\0&1&2\end{pmatrix},
  \qquad
  \vb=\begin{pmatrix}1\\0\\0\end{pmatrix}.
\]
$\mat{A}$ is symmetric, and its eigenvalues are $2-\sqrt2,\,2,\,2+\sqrt2$, all
positive, so it is SPD---the setting CG (\cref{ch:cg}) will live in.

\textbf{Build the Krylov sequence.} Each new vector is one operator apply on the
last. Starting from $\vb$:
\[
  \vb=\begin{pmatrix}1\\0\\0\end{pmatrix},\qquad
  \mat{A}\vb=\begin{pmatrix}2\\1\\0\end{pmatrix},\qquad
  \mat{A}^2\vb=\mat{A}(\mat{A}\vb)=\begin{pmatrix}5\\4\\1\end{pmatrix}.
\]
(Check the last: $\mat{A}(2,1,0)^{\mathsf T}=(2\!\cdot\!2+1,\ 2+2,\ 1)^{\mathsf
T}=(5,4,1)^{\mathsf T}$.) Two operator applies built the three spanning vectors. The
three Krylov subspaces are
\[
  \mathcal{K}_1=\operatorname{span}\{\vb\},\quad
  \mathcal{K}_2=\operatorname{span}\{\vb,\mat{A}\vb\},\quad
  \mathcal{K}_3=\operatorname{span}\{\vb,\mat{A}\vb,\mat{A}^2\vb\}.
\]

\tcblower
\textbf{Are the three vectors independent?} Stack them as columns and take the
determinant:
\[
  \det\!\begin{pmatrix}1&2&5\\0&1&4\\0&0&1\end{pmatrix}=1\neq0 .
\]
The matrix is upper-triangular with unit diagonal, so the three vectors are linearly
independent: $\mathcal{K}_3=\Real^3$, the whole space. The nesting
$\mathcal{K}_1\subsetneq\mathcal{K}_2\subsetneq\mathcal{K}_3$ is strict, one new
dimension per apply.

\textbf{The solution lives in $\mathcal{K}_3$.} Solve $\mat{A}\vx=\vb$ directly for the
target. With $\det\mat{A}=2(4-1)-1(2-0)=4$, the solution is
\[
  \vx^\star=\mat{A}^{-1}\vb=\tfrac14\begin{pmatrix}3\\-2\\1\end{pmatrix}
   =\begin{pmatrix}0.75\\-0.5\\0.25\end{pmatrix}.
\]
Since $\mathcal{K}_3=\Real^3$, $\vx^\star$ is certainly in $\mathcal{K}_3$---and we can
exhibit it as a polynomial in $\mat{A}$ applied to $\vb$, the form
\cref{sec:krylov-poly} promised. Seek $\vx^\star=c_0\vb+c_1\mat{A}\vb+c_2\mat{A}^2\vb$:
\[
  c_0\begin{pmatrix}1\\0\\0\end{pmatrix}
  +c_1\begin{pmatrix}2\\1\\0\end{pmatrix}
  +c_2\begin{pmatrix}5\\4\\1\end{pmatrix}
  =\begin{pmatrix}0.75\\-0.5\\0.25\end{pmatrix}.
\]
The bottom row gives $c_2=0.25$; the middle row $c_1+4c_2=-0.5$ gives $c_1=-1.5$; the
top row $c_0+2c_1+5c_2=0.75$ gives $c_0=0.75+3-1.25=2.5$. So
\[
  \vx^\star=\bigl(2.5\,I-1.5\,\mat{A}+0.25\,\mat{A}^2\bigr)\vb
  = q(\mat{A})\,\vb,\quad q(t)=0.25t^2-1.5t+2.5 .
\]
The exact inverse-times-$\vb$ is a degree-2 polynomial in $\mat{A}$ applied to
$\vb$---exactly the minimal-polynomial mechanism of \cref{sec:krylov-poly}, here with
$d=3$ (so degree $d-1=2$). A Krylov method on this system reaches the exact answer in
at most $3$ steps; CG, minimizing the error over each $\mathcal{K}_m$, would land on
$\vx^\star$ by step $3$ in exact arithmetic---and, because the spectrum is tightly
clustered around $2$, would be very close well before that. The subspace built from
nothing but operator applies contains the solution, and we found the coefficients
that pick it out.
\end{workedexample}

\begin{workedexample}{Costing one Krylov step against a direct factorization}{krylovcost}
We make the cost invariant concrete by pricing one step of a Krylov method and
comparing it to a sparse direct factorization, on a system the size a real solve
faces. Take a 3-D finite-element operator with $N=10^6$ unknowns and a bounded
$\nu=27$ nonzeros per row (a typical low-order stencil), so $\mat{A}$ has about
$\nu N=2.7\times10^7$ nonzeros.

\textbf{Cost of one Krylov step.} The kernel of \cref{sec:krylov-kernel} does, per
step:
\begin{itemize}[nosep]
  \item \emph{one operator apply} $\vw=\mat{A}\vp$. A sparse apply touches each
  nonzero once: $\approx\nu N=2.7\times10^7$ multiply--adds. This is the dominant
  cost.
  \item \emph{a couple of inner products / norms}, e.g.\ $\inner{\vp}{\vw}$ and
  $\norm{\vr}$: each $N$ multiply--adds, so $\approx 2N=2\times10^6$.
  \item \emph{a couple of \emph{axpy} updates}, e.g.\ $\vx\leftarrow\vx+\alpha\vp$ and
  $\vr\leftarrow\vr-\alpha\vw$: each $N$ multiply--adds, so $\approx 2N=2\times10^6$.
\end{itemize}
Total: about $\nu N+4N=(27+4)\times10^6\approx 3.1\times10^7$ operations, and the apply
is $\approx87\%$ of it. The step is $O(N)$, and it is one apply plus small change.

\tcblower
\textbf{Cost of a sparse direct factorization.} A sparse Cholesky on a 3-D mesh costs
about $O(N^2)$ arithmetic, with the fill-in of \cref{sec:krylov-direct} producing
factors of about $O(N^{4/3})$ nonzeros. Even taking a generous constant, the
factorization work at $N=10^6$ is on the order of $10^{12}$ operations or more---and
the factors need on the order of $N^{4/3}=10^8$ stored nonzeros, several times the
$2.7\times10^7$ of $\mat{A}$ itself, before they are even applied.

\textbf{The comparison.} Put the two side by side:
\[
  \underbrace{3.1\times10^7}_{\text{one Krylov step}}
  \qquad\text{versus}\qquad
  \underbrace{\sim 10^{12}}_{\text{one direct factorization}} .
\]
The factorization costs as much as roughly $10^{12}/(3.1\times10^7)\approx3\times10^4$
Krylov steps. So a preconditioned Krylov solve that converges in, say, a few hundred
steps wins by orders of magnitude---\emph{and} never forms a matrix, since each of
those steps needs only the matrix-free apply (\cref{ch:matrixfree}), while the
factorization needs the entries and pays the fill-in to store its denser factors.

\textbf{The reading.} The whole economics of large-scale solving is in this
comparison. A direct method pays a huge one-time $O(N^2)$ cost and then solves
cheaply; an iterative method pays a small $O(N)$ cost per step and bets that a
bounded number of steps suffice. For a million-unknown 3-D system, the bet pays:
the per-step cost is dominated by one sparse apply, the apply can be matrix-free,
and a good preconditioner (\cref{ch:precond}) keeps the step count bounded. ``One
apply per step,'' costed out, is why Part~IV exists.
\end{workedexample}

\section{The map of Part~IV}
\label{sec:krylov-map}

This chapter laid the shared foundation; the rest of Part~IV fills in the
specializations and the machinery that makes them fast. Here is the road ahead, each
piece a refinement of something built in this chapter.

\begin{description}[style=nextline, leftmargin=1.4em]
  \item[Conjugate gradients (\cref{ch:cg}).] The Krylov method for SPD systems---
  electrostatics, magnetostatics, implicit time steps. It specializes
  \code{krylov\_step} with a short two-term recurrence (\cref{ch:fold}'s
  \code{iterate\_while\_with\_prev}) and no explicit orthogonalization, minimizing
  the error in the $\mat{A}$-norm. One apply per step.
  \item[GMRES and FGMRES (\cref{ch:gmres}).] The Krylov method for general
  (nonsymmetric, complex) systems---the driven analysis. It builds an explicit
  orthonormal Krylov basis and minimizes the residual norm over it, which forces the
  orthogonalization stage and an outer restart loop. FGMRES threads a second basis to
  carry a per-step (flexible) preconditioner, exactly the per-step variant that
  \cref{ch:operators}'s decision rule routes into the state.
  \item[Orthogonalization (\cref{ch:orthog}).] The stage GMRES needs: making each new
  basis vector orthogonal to the previous ones, and the modified-versus-classical
  Gram--Schmidt choice that is the orthogonalization-variant non-law of
  \cref{sec:krylov-laws}.
  \item[The least-squares core (\cref{ch:givens}).] How GMRES actually picks the
  minimizing iterate inside its basis: a small dense least-squares problem solved
  incrementally with Givens rotations as the basis grows.
  \item[Preconditioning and multigrid (\cref{ch:precond,ch:multigrid}).] The
  conditioning half of the story---replacing $\mat{A}$ with an equivalent
  better-conditioned system so the spectrum clusters (\cref{fig:krylov-convergence})
  and the step count drops. The preconditioner is itself a constructed operator
  (\cref{ch:operators}) the kernel applies; multigrid is the most powerful one, an
  approximate inverse built from a hierarchy of meshes.
\end{description}

Every one of these is a setting of the kernel and the fold built here, or a piece of
machinery that feeds them. Learn the Krylov subspace, the one-apply-per-step kernel,
and the solve-as-fold-as-operator reading, and the rest of Part~IV is
specialization.

\summarysection
Discretization (Part~II) produces a large \emph{sparse} linear system
$\mat{A}\vx=\vb$: $\mat{A}$ has $N$ up to the tens of millions but only $O(N)$
nonzeros, because the finite-element basis has local support and locality is
sparsity. \emph{Direct} methods (Gaussian elimination, $LU$/Cholesky) are too
expensive at scale---fill-in makes the factors far denser than $\mat{A}$, costing
about $O(N^2)$ in 3-D and $O(N^3)$ dense---and they need the matrix entries, which a
matrix-free operator does not provide. \emph{Iterative} methods win because they need
only the operator's \emph{action} $\vv\mapsto\mat{A}\vv$, the exact interface
\cref{ch:matrixfree} supplies, and cost close to $O(N)$ when well preconditioned. A
\emph{Krylov subspace} $\mathcal{K}_m(\mat{A},\vb)=\operatorname{span}\{\vb,\mat{A}\vb,
\dots,\mat{A}^{m-1}\vb\}$ is built from nothing but repeated applies; its nested
ladder contains better and better approximations, and the minimal polynomial puts the
exact solution in $\mathcal{K}_d$ for some $d\le N$ (so Krylov terminates in $\le N$
steps in exact arithmetic). Convergence speed is governed by the spectrum: a
\emph{clustered} spectrum converges fast, a \emph{spread} one slowly---which is why
preconditioning exists. Every Krylov method shares one \emph{step kernel}
\code{krylov\_step}: one operator apply, a few inner products and norms, a few
\emph{axpy} updates, an optional orthogonalization, a state update---with the cost
invariant of \emph{exactly one apply per step}. The outer loop is the
\code{iterate\_while} fold of \cref{ch:fold}, and the whole solve is a
\emph{constructed operator} (\cref{ch:operators}), \code{ksp\_solve}: an approximate
inverse you build once and apply, obeying linearity and the near-inverse law modulo
tolerance, and deliberately \emph{not} obeying bit-determinism across reduction
order, orthogonalization, or initial guess---a controlled approximation whose
accuracy is a tolerance and whose speed is conditioning. Part~IV specializes this
foundation: CG (\cref{ch:cg}), GMRES/FGMRES (\cref{ch:gmres}), orthogonalization
(\cref{ch:orthog}), the least-squares core (\cref{ch:givens}), and preconditioning
and multigrid (\cref{ch:precond,ch:multigrid}).

\furtherreading
The standard reference for iterative methods on sparse systems is
Saad~\citep{saad2003}; its early chapters develop the Krylov subspace, the
polynomial view, and the spectrum-driven convergence bounds at exactly the level of
this chapter, and the later chapters are the canonical source for the CG and GMRES
specializations of \cref{ch:cg,ch:gmres} and for preconditioning
(\cref{ch:precond}). For the numerical-linear-algebra foundations---why fill-in
makes direct methods scale poorly, the conditioning and floating-point subtleties
behind the non-laws of \cref{sec:krylov-laws}---Golub and Van
Loan~\citep{golubvanloan2013} is the comprehensive reference. Trefethen and
Bau~\citep{trefethenbau1997} give the most readable account of the
polynomial-approximation picture of Krylov convergence (their lectures on the
Arnoldi and CG processes are especially clear on why a clustered spectrum converges
fast), and are an ideal companion to this chapter's
\cref{sec:krylov-poly,sec:krylov-convergence}.

\exercisessection
\begin{enumerate}[nosep]
  \item \textbf{Sparsity counts.} A 3-D finite-element operator has $N=4\times10^6$
  unknowns and an average of $\nu=30$ nonzeros per row. (a) How many numbers does the
  sparse operator store? (b) How many would the dense $N\times N$ array hold? (c) By
  what factor does sparsity reduce the storage? (d) In one sentence, tie the bounded
  $\nu$ back to the local support of the basis functions (\cref{ch:refelement}).
  \item \textbf{Why not factor.} Using \cref{fig:krylov-cost}, estimate the ratio of
  work between a dense direct solve ($O(N^3)$) and a near-linear iterative solve
  ($\approx O(N)$) at $N=10^5$ and again at $N=10^7$. By what factor does the
  iterative advantage grow as $N$ goes up by $100\times$? Explain in terms of the
  exponents why the gap widens.
  \item \textbf{Build a Krylov subspace.} For
  $\mat{A}=\bigl(\begin{smallmatrix}3&1\\1&3\end{smallmatrix}\bigr)$ and
  $\vb=(1,1)^{\mathsf T}$: (a) compute $\mat{A}\vb$; (b) is $\mathcal{K}_2(\mat{A},\vb)
  =\operatorname{span}\{\vb,\mat{A}\vb\}$ all of $\Real^2$, or does it collapse?
  (c) Solve $\mat{A}\vx=\vb$ directly and find scalars $c_0,c_1$ with
  $\vx=c_0\vb+c_1\mat{A}\vb$, exhibiting $\vx=q(\mat{A})\vb$. (d) Why does the collapse
  in (b) mean a Krylov method finishes in \emph{one} step here? (Hint: $\vb$ is an
  eigenvector.)
  \item \textbf{The polynomial view.} \Cref{sec:krylov-poly} shows
  $\mat{A}^{-1}=q(\mat{A})$ for a polynomial $q$ of degree $d-1$, where $d$ is the
  degree of the minimal polynomial. If $\mat{A}$ has only $3$ distinct eigenvalues
  (each possibly repeated), what is the most steps a Krylov method that finds the best
  iterate in each $\mathcal{K}_m$ can need, in exact arithmetic, regardless of $N$?
  Why does the \emph{number of distinct eigenvalues}, not $N$, set the bound?
  \item \textbf{Conditioning and speed.} Two SPD systems have condition numbers
  $\kappa_1=100$ and $\kappa_2=10^4$. Using the CG rule of thumb that the iteration
  count scales like $\sqrt{\kappa}$, estimate the ratio of step counts to reach a
  fixed tolerance. If a preconditioner clusters the spectrum of the second system so
  its effective condition number drops to $100$, what happens to its step count?
  Relate your answer to \cref{fig:krylov-convergence}.
  \item \textbf{One apply per step.} Reread \cref{wex:krylovcost}. (a) For
  $N=10^6$ and $\nu=27$, what fraction of one Krylov step's arithmetic is the operator
  apply? (b) Suppose a colleague proposes computing \emph{two} extra inner products per
  step for a fancier convergence test. By what percentage does that raise the per-step
  cost, and why is it nonetheless usually a bad trade if it does not reduce the step
  count? (c) Why does the cost invariant ``one apply per step'' survive unchanged when
  the apply is matrix-free (\cref{ch:matrixfree})?
  \item \textbf{The solve is a fold.} Write the convergence predicate of a Krylov
  solve that stops when either the relative residual drops below $10^{-8}$ or the
  iteration count reaches $500$, as a pure function of the carry (cf.\
  \cref{sec:krylov-fold} and the predicate rule of \cref{sec:predicaterule}). Which
  fields must the carry contain for this predicate to be expressible, and why can the
  predicate \emph{not} read the residual the current step just produced as an extra?
  \item \textbf{The solve is an operator.} The driven analysis (\cref{ch:situating})
  rebuilds its operator $\mat{A}(\omega)$ at every frequency, while electrostatics
  reuses one operator across many right-hand sides. Using
  \cref{fig:krylov-ksp} and the constructed-operator pattern of \cref{ch:operators},
  explain which analysis can hoist the \code{ksp\_solve} construction out of its loop
  and which cannot, and why the difference is one of \emph{kind}, not degree.
  \item \textbf{The non-laws.} A caller writes an iterative-refinement loop that
  assumes $\mat{A}\,\code{ksp\_solve}(K,\vb).\vx=\vb$ \emph{exactly} after a solve.
  (a) Which non-law of \cref{sec:krylov-laws} does this violate? (b) What goes wrong,
  concretely, if the assumption is false by a quantity near the tolerance? (c) Name
  the three bit-level non-determinism sources a Krylov solve carries, and for each,
  state whether it changes the \emph{mathematical} solution or only its
  floating-point realization.
\end{enumerate}
